\section{Orbits, Cycles, and the Alternating Groups}
\textbf{Conventions.} Permutations compose from right to left. The identity is $e$. A group generated by an element $a$ is denoted $\gen{a}$.

\exercise{9.1}{Orbits of a finite permutation}
\idea Start at an unused symbol, repeatedly apply the permutation until it returns, and repeat with a symbol not yet used.
The nontrivial cycles are $(1,5,2)(4,6)$. Each fixed point forms a singleton orbit, so the complete orbit list is\[\{1,2,5\},\quad \{4,6\},\quad \{3\}.\]Every symbol occurs in exactly one listed set; negative powers move backwards within the same set and produce no further points.

\exercise{9.2}{Orbits of a finite permutation}
\idea Start at an unused symbol, repeatedly apply the permutation until it returns, and repeat with a symbol not yet used.
The nontrivial cycles are $(1,5,8,7)(2,6,3)$. Each fixed point forms a singleton orbit, so the complete orbit list is\[\{1,5,7,8\},\quad \{2,3,6\},\quad \{4\}.\]Every symbol occurs in exactly one listed set; negative powers move backwards within the same set and produce no further points.

\exercise{9.3}{Orbits of a finite permutation}
\idea Start at an unused symbol, repeatedly apply the permutation until it returns, and repeat with a symbol not yet used.
The nontrivial cycles are $(1,2,3,5,4)(7,8)$. Each fixed point forms a singleton orbit, so the complete orbit list is\[\{1,2,3,4,5\},\quad \{7,8\},\quad \{6\}.\]Every symbol occurs in exactly one listed set; negative powers move backwards within the same set and produce no further points.

\exercise{9.4}{Orbits of the integer shift by $1$}
The $r$th power sends an integer $a$ to $a+1r$, where $r$ ranges over $\Z$. Hence the orbit of $a$ is $a+1\Z$, its residue class modulo $1$. The distinct orbits are\[\Z.\]The sign of the shift only changes the direction in which these same points are traversed.

\exercise{9.5}{Orbits of the integer shift by $2$}
The $r$th power sends an integer $a$ to $a+2r$, where $r$ ranges over $\Z$. Hence the orbit of $a$ is $a+2\Z$, its residue class modulo $2$. The distinct orbits are\[0+2\Z,\quad 1+2\Z.\]The sign of the shift only changes the direction in which these same points are traversed.

\exercise{9.6}{Orbits of the integer shift by $-3$}
The $r$th power sends an integer $a$ to $a+-3r$, where $r$ ranges over $\Z$. Hence the orbit of $a$ is $a+3\Z$, its residue class modulo $3$. The distinct orbits are\[0+3\Z,\quad 1+3\Z,\quad 2+3\Z.\]The sign of the shift only changes the direction in which these same points are traversed.

\exercise{9.7}{A product of cycles}
\idea Read the factors from right to left, leaving a symbol unchanged whenever it is absent from a factor.
The images of the eight arguments are\[\begin{pmatrix}1 & 2 & 3 & 4 & 5 & 6 & 7 & 8\\4 & 1 & 3 & 5 & 8 & 6 & 2 & 7\end{pmatrix}.\]Following each chain of these images yields the disjoint-cycle answer\[(1,4,5)(7,8)(2,5,7)=(1,4,5,8,7,2).\]Any symbol omitted from the final cycles is fixed. The two-row form provides a check of every image, not just of one cycle.

\exercise{9.8}{A product of cycles}
\idea Read the factors from right to left, leaving a symbol unchanged whenever it is absent from a factor.
The images of the eight arguments are\[\begin{pmatrix}1 & 2 & 3 & 4 & 5 & 6 & 7 & 8\\3 & 7 & 2 & 8 & 5 & 4 & 1 & 6\end{pmatrix}.\]Following each chain of these images yields the disjoint-cycle answer\[(1,3,2,7)(4,8,6)=(1,3,2,7)(4,8,6).\]Any symbol omitted from the final cycles is fixed. The two-row form provides a check of every image, not just of one cycle.

\exercise{9.9}{A product of cycles}
\idea Read the factors from right to left, leaving a symbol unchanged whenever it is absent from a factor.
The images of the eight arguments are\[\begin{pmatrix}1 & 2 & 3 & 4 & 5 & 6 & 7 & 8\\5 & 4 & 3 & 7 & 8 & 6 & 2 & 1\end{pmatrix}.\]Following each chain of these images yields the disjoint-cycle answer\[(1,2)(4,7,8)(2,1)(7,2,8,1,5)=(1,5,8)(2,4,7).\]Any symbol omitted from the final cycles is fixed. The two-row form provides a check of every image, not just of one cycle.

\exercise{9.10}{Disjoint cycles and transpositions}
\step{1. Follow the images.} The nontrivial chains are\[1\mapsto8\mapsto1;\quad 3\mapsto6\mapsto4\mapsto3;\quad 5\mapsto7\mapsto5.\]Thus the disjoint-cycle expression is $(1,8)(3,6,4)(5,7)$; any missing point is fixed.
\step{2. Split each cycle.} Use $(a_1,a_2,\ldots,a_k)=(a_1,a_k)\cdots(a_1,a_3)(a_1,a_2)$, verified by following each $a_i$. This gives\[(1,8)(3,6,4)(5,7)=(1,8)(3,4)(3,6)(5,7).\]

\exercise{9.11}{Disjoint cycles and transpositions}
\step{1. Follow the images.} The nontrivial chains are\[1\mapsto3\mapsto4\mapsto1;\quad 2\mapsto6\mapsto2;\quad 5\mapsto8\mapsto7\mapsto5.\]Thus the disjoint-cycle expression is $(1,3,4)(2,6)(5,8,7)$; any missing point is fixed.
\step{2. Split each cycle.} Use $(a_1,a_2,\ldots,a_k)=(a_1,a_k)\cdots(a_1,a_3)(a_1,a_2)$, verified by following each $a_i$. This gives\[(1,3,4)(2,6)(5,8,7)=(1,4)(1,3)(2,6)(5,7)(5,8).\]

\exercise{9.12}{Disjoint cycles and transpositions}
\step{1. Follow the images.} The nontrivial chains are\[1\mapsto3\mapsto4\mapsto7\mapsto8\mapsto6\mapsto5\mapsto2\mapsto1.\]Thus the disjoint-cycle expression is $(1,3,4,7,8,6,5,2)$; any missing point is fixed.
\step{2. Split each cycle.} Use $(a_1,a_2,\ldots,a_k)=(a_1,a_k)\cdots(a_1,a_3)(a_1,a_2)$, verified by following each $a_i$. This gives\[(1,3,4,7,8,6,5,2)=(1,2)(1,5)(1,6)(1,8)(1,7)(1,4)(1,3).\]

\exercise{9.13}{Orders of cycles and permutations}
\begin{enumerate}[label=(\alph*)]
\item $(1,4,5,7)$ has order $4$: four applications return each listed point, and no smaller positive number returns $1$.
\item A cycle of length $k$ has order $k$. Each application advances one step, so a power fixes the cycle exactly when the exponent is divisible by $k$.
\item The disjoint cycles of $\sigma=(4,5)(2,3,7)$ have lengths $2,3$, giving order $6$. Those of $\tau=(1,4)(3,5,7,8)$ have lengths $2,4$, giving order $4$.
\item Exercises 10, 11, 12 have nontrivial cycle lengths $(2,3,2)$, $(3,2,3)$, and $(8)$, respectively. Their orders are $6,6,8$.
\item The order of a permutation is the least common multiple of its disjoint-cycle lengths: a power is the identity precisely when it is a multiple of every cycle length.
\end{enumerate}

\exercise{9.14}{Maximum element order in $S_{5}$}
\idea Disjoint-cycle lengths are positive integers summing to $5$, and the order is their least common multiple. Thus check partitions, rather than individual permutations.
Grouping the partitions by their largest part $L$, the largest possible least common multiple in each group is\[\begin{array}{c|rrrrr}L&1&2&3&4&5\\\hline \text{maximum}&1&2&6&4&5\end{array}\]For example, these values can be reproduced by listing the remaining nondecreasing parts, each at most $L$, with sum $5-L$, and taking the least common multiple with $L$. The overall maximum is $6$, attained by disjoint cycles of lengths $2,3$. This gives both an example and an exhaustive upper bound.

\exercise{9.15}{Maximum element order in $S_{6}$}
\idea Disjoint-cycle lengths are positive integers summing to $6$, and the order is their least common multiple. Thus check partitions, rather than individual permutations.
Grouping the partitions by their largest part $L$, the largest possible least common multiple in each group is\[\begin{array}{c|rrrrrr}L&1&2&3&4&5&6\\\hline \text{maximum}&1&2&6&4&5&6\end{array}\]For example, these values can be reproduced by listing the remaining nondecreasing parts, each at most $L$, with sum $6-L$, and taking the least common multiple with $L$. The overall maximum is $6$, attained by disjoint cycles of lengths $1,2,3$. This gives both an example and an exhaustive upper bound.

\exercise{9.16}{Maximum element order in $S_{7}$}
\idea Disjoint-cycle lengths are positive integers summing to $7$, and the order is their least common multiple. Thus check partitions, rather than individual permutations.
Grouping the partitions by their largest part $L$, the largest possible least common multiple in each group is\[\begin{array}{c|rrrrrrr}L&1&2&3&4&5&6&7\\\hline \text{maximum}&1&2&6&12&10&6&7\end{array}\]For example, these values can be reproduced by listing the remaining nondecreasing parts, each at most $L$, with sum $7-L$, and taking the least common multiple with $L$. The overall maximum is $12$, attained by disjoint cycles of lengths $3,4$. This gives both an example and an exhaustive upper bound.

\exercise{9.17}{Maximum element order in $S_{10}$}
\idea Disjoint-cycle lengths are positive integers summing to $10$, and the order is their least common multiple. Thus check partitions, rather than individual permutations.
Grouping the partitions by their largest part $L$, the largest possible least common multiple in each group is\[\begin{array}{c|rrrrrrrrrr}L&1&2&3&4&5&6&7&8&9&10\\\hline \text{maximum}&1&2&6&12&30&12&21&8&9&10\end{array}\]For example, these values can be reproduced by listing the remaining nondecreasing parts, each at most $L$, with sum $10-L$, and taking the least common multiple with $L$. The overall maximum is $30$, attained by disjoint cycles of lengths $2,3,5$. This gives both an example and an exhaustive upper bound.

\exercise{9.18}{Maximum element order in $S_{15}$}
\idea Disjoint-cycle lengths are positive integers summing to $15$, and the order is their least common multiple. Thus check partitions, rather than individual permutations.
Grouping the partitions by their largest part $L$, the largest possible least common multiple in each group is\[\begin{array}{c|rrrrrrrrrrrrrrr}L&1&2&3&4&5&6&7&8&9&10&11&12&13&14&15\\\hline \text{maximum}&1&2&6&12&60&60&105&56&45&30&44&12&26&14&15\end{array}\]For example, these values can be reproduced by listing the remaining nondecreasing parts, each at most $L$, with sum $15-L$, and taking the least common multiple with $L$. The overall maximum is $105$, attained by disjoint cycles of lengths $3,5,7$. This gives both an example and an exhaustive upper bound.

\exercise{9.19}{Labeling the Cayley digraph of $A_4$}
Put $a=(1,2,3)$ and $b=(1,2)(3,4)$. Follow each solid arrow by right multiplying by $a$, and each dashed edge by right multiplying by $b$. The following table labels every position in the supplied diagram; the letters also appear in the schematic below.
\[\begin{array}{c|l|l}
\text{vertex}&\text{position}&\text{permutation}\\\hline
A & top left & e\\
B & top right & (1,2,3)\\
C & below top pair & (1,3,2)\\
D & left outside & (1,2)(3,4)\\
E & left inner & (2,4,3)\\
F & bottom left & (1,4,3)\\
G & middle top & (2,3,4)\\
H & middle right & (1,3)(2,4)\\
I & middle left & (1,4,2)\\
J & right outside & (1,3,4)\\
K & bottom right & (1,2,4)\\
L & right inner & (1,4)(2,3)\\
\end{array}\]
\begin{center}\setlength{\unitlength}{1mm}\begin{picture}(100,88)
\put(39,83){$A$}\put(61,83){$B$}\put(50,65){$C$}
\put(10,28){$D$}\put(29,31){$E$}\put(20,10){$F$}
\put(50,51){$G$}\put(40,35){$I$}\put(60,35){$H$}
\put(90,28){$J$}\put(80,10){$K$}\put(70,31){$L$}
\put(43,82){\vector(1,0){18}}\put(62,79){\vector(-3,-5){10}}\put(49,68){\vector(-3,5){9}}
\put(13,27){\vector(6,1){16}}\put(29,27){\vector(-2,-5){6}}\put(19,14){\vector(-3,5){8}}
\put(43,39){\vector(3,5){8}}\put(53,50){\vector(3,-5){8}}\put(60,34){\vector(-1,0){16}}
\put(73,30){\vector(6,-1){16}}\put(91,25){\vector(-3,-5){8}}\put(79,14){\vector(-2,5){6}}
\put(12,31){\line(1,2){26}}\put(65,79){\line(1,-2){24}}
\put(52,63){\line(0,-1){9}}\put(32,32){\line(1,0){8}}\put(64,34){\line(1,0){6}}\put(24,12){\line(1,0){54}}
\end{picture}\end{center}
The six edges without arrowheads represent the dashed $b$ edges; both directions give the same operation since $b^2=e$. For instance $D=b$, $E=ba$, and $F=ba^2$. The twelve entries are distinct and exhaust $A_4$, checking that no vertex has been missed.

\exercise{9.20}{Definition of an orbit}
The proposed definition is correct when ``minimal'' means minimal under inclusion among nonempty invariant subsets: an orbit is a nonempty subset $O\subseteq A$ with $\sigma[O]=O$ and no proper nonempty subset having that property. Indeed, $O_a=\{\sigma^k(a):k\in\Z\}$ is invariant, and any nonempty invariant subset containing a point of $O_a$ contains its entire two-sided orbit. Conversely, pick $a$ in a minimal such set; its orbit lies inside the set and hence must equal it.

\exercise{9.21}{Definition of a cycle}
A cycle in $S_n$ is a permutation with \emph{at most one orbit having more than one element}. Its length is the size of its largest orbit. Fixed points outside the moving cycle produce additional singleton orbits, so ``only one orbit'' is too restrictive. The book includes the identity as a cycle of length $1$.

\exercise{9.22}{Definition of the alternating group}
Specify the common finite underlying set: the alternating group $A_n$ is the subgroup of $S_n$ consisting of all even permutations of $\{1,\ldots,n\}$. Without specifying this set or $n$, ``all even permutations'' does not describe a single group. Closure follows because products and inverses of even permutations are even.

\exercise{9.23}{True or false}
\begin{enumerate}[label=(\alph*)]
\item False: $(1,2)(3,4)$ has two nontrivial orbits.
\item True: a cycle is a special kind of permutation.
\item False: first one must prove that the parity of the number of transpositions is independent of the chosen expression; otherwise the definition could assign both parities to one permutation.
\item False: $\gen{(1,2,3,4)}$ has an odd permutation but its elements are $e$, a $4$-cycle, a double transposition, and the inverse $4$-cycle; none is a transposition.
\item False: $|A_5|=5!/2=60$.
\item False: $S_2=\{e,(1,2)\}$ is cyclic.
\item True: $A_3$ is cyclic of order $3$.
\item True: extend each permutation of $\{1,\ldots,7\}$ by fixing $8$.
\item True: relabel the seven symbols other than $5$, and extend by fixing $5$; the particular omitted symbol is immaterial.
\item False: the identity is even and so is absent, and the product of two odd permutations is even.
\end{enumerate}

\exercise{9.24}{The table of $A_3$}
The even permutations are $\rho_0=e$, $\rho_1=(1,2,3)$, and $\rho_2=(1,3,2)$. Each $3$-cycle is a product of two transpositions, while $\mu_1,\mu_2,\mu_3$ are single transpositions and are odd. Thus
\[\begin{array}{c|ccc}&\rho_0&\rho_1&\rho_2\\\hline
\rho_0&\rho_0&\rho_1&\rho_2\\
\rho_1&\rho_1&\rho_2&\rho_0\\
\rho_2&\rho_2&\rho_0&\rho_1
\end{array}\]
The entries follow from $\rho_1^3=e$ and $\rho_2=\rho_1^2$.

\exercise{9.25}{Synopsis of the determinant proof}
A permutation matrix obtained from the identity by $k$ row transpositions has determinant $(-1)^k$, so two transposition expressions for the same permutation must have the same parity because one matrix cannot have determinants $1$ and $-1$ simultaneously.

\exercise{9.26}{Synopsis of the orbit-counting proof}
Multiplying a permutation by one transposition merges two distinct orbits or splits one orbit in two, changing its number of orbits by exactly $1$. Starting from the identity's $n$ orbits, the parity of the number of transpositions is therefore determined by the final number of orbits and cannot depend on the expression.

\exercise{9.27}{Bounds on numbers of transpositions}
Let the disjoint-cycle lengths, including fixed points, be $k_1,\ldots,k_t$, with sum $n$. A $k_i$-cycle uses $k_i-1$ transpositions, so the entire permutation uses
\[(k_1-1)+\cdots+(k_t-1)=n-t.\]
\textbf{(a)} Since $t\ge1$, this is at most $n-1$. The identity can be represented by zero transpositions, the empty product.
\textbf{(b)} A permutation which is not a cycle has at least two nontrivial orbits, hence $t\ge2$, giving at most $n-2$ transpositions.
\textbf{(c)} If the permutation is odd, the displayed number $m=n-t$ is odd. Since $m\le n-1<2n+3$ and $2n+3-m$ is even, insert $(2n+3-m)/2$ pairs $(1,2)(1,2)$ into that expression to obtain exactly $2n+3$ factors. For an even permutation $m$ is even and the same insertion of $(2n+8-m)/2$ pairs gives exactly $2n+8$ factors. Each pair is the identity and so does not change the permutation.

\exercise{9.28}{Merging orbits when a point is fixed}
\textbf{(a)} Write the orbit containing $j$ as $(j,x_1,\ldots,x_k)$ and suppose $i$ is a different fixed point. Then direct calculation gives
\[(i,j)(j,x_1,\ldots,x_k)=(i,j,x_1,\ldots,x_k).\]
Thus the two old orbit diagrams and the new diagram can be drawn as
\[\underbrace{i\longrightarrow i}_{\text{one orbit}}\qquad
\underbrace{j\longrightarrow x_1\longrightarrow\cdots\longrightarrow x_k\longrightarrow j}_{\text{one orbit}}
\quad\rightsquigarrow\quad
 i\longrightarrow j\longrightarrow x_1\longrightarrow\cdots\longrightarrow x_k\longrightarrow i.
\]
The other orbits are unchanged, so the orbit count decreases by one.
\textbf{(b)} If $j$ is fixed too, the two singleton loops $i\to i$ and $j\to j$ become the single $2$-cycle $i\to j\to i$ under $(i,j)\sigma$. Again exactly two orbits have become one.

\exercise{9.29}{Even permutations in a subgroup}
If $H$ has no odd permutations, every element is even. Otherwise fix an odd $\tau\in H$. Left multiplication by $\tau$ maps the even elements of $H$ to odd elements of $H$, because signs multiply. Its inverse is left multiplication by $\tau^{-1}$, which is also odd. Thus it is a bijection between the even and odd elements of $H$. As $H$ is finite, exactly half its elements are even.

\exercise{9.30}{How many points a cycle moves}
A cycle of length $n>1$ moves exactly its $n$ listed points: each goes to the next distinct listed point, and all other points are fixed. Under the book's convention, a cycle of length $1$ is the identity and moves $0$ points. Thus the usual answer $n$ requires the nontrivial-cycle condition $n>1$.

\exercise{9.31}{Finitely supported permutations form a subgroup}
For a permutation $\sigma$, write $\operatorname{supp}(\sigma)=\{a:\sigma(a)\ne a\}$. The identity has empty support. If both $\sigma$ and $\tau$ fix a point $a$, their product fixes $a$, so
\[\operatorname{supp}(\sigma\tau)\subseteq\operatorname{supp}(\sigma)\cup\operatorname{supp}(\tau).\]
A union of two finite sets is finite, proving product closure. A point is fixed by $\sigma$ if and only if it is fixed by $\sigma^{-1}$, so the two supports are equal. Hence inverses also have finite support, and the given set is a subgroup.

\exercise{9.32}{A uniform bound on support size}
No. Choose $100$ distinct points of the infinite set $A$, and let $\sigma$ be a $50$-cycle on the first $50$, while $\tau$ is a $50$-cycle on the other $50$. Extend both by the identity elsewhere. Each moves exactly $50$ points, so both belong to $K$. Their disjoint supports mean $\sigma\tau$ moves all $100$ selected points, so it is not in $K$. Thus product closure fails.

\exercise{9.33}{The odd permutations form a coset}
Let $\sigma$ be the fixed odd permutation and let $\pi$ be any odd permutation. Set $\alpha=\sigma^{-1}\pi$. Both factors are odd, so their product is even and $\alpha\in A_n$. Then $\pi=\sigma\alpha$. Conversely, the product of $\sigma$ with any even permutation is odd. Thus the set of odd permutations is exactly $\sigma A_n$.

\exercise{9.34}{Squaring an odd-length cycle}
Let $\sigma=(a_0,a_1,\ldots,a_{n-1})$ with $n$ odd. Its square moves $a_i$ to $a_{i+2\bmod n}$. If $2k\equiv0\pmod n$, oddness implies $n\mid k$, so the first positive return occurs after $n$ steps. Hence all $n$ symbols lie in one cycle of $\sigma^2$, and every point outside the original cycle stays fixed. For $n=1$, both permutations are the identity, also a cycle by convention.

\exercise{9.35}{When a power of a cycle is a cycle}
\idea Advancing $r$ steps around an $n$-cycle partitions its $n$ points into $d=\gcd(n,r)$ cycles, each of length $n/d$.
Indeed, $\sigma^{rk}(a)=a$ precisely when $n\mid rk$, whose least positive solution is $k=n/d$. If $n/d>1$, all $d$ orbits are nontrivial, so the power is a cycle exactly when $d=1$.
There is an important convention in this text: the identity is also a cycle. If $n\mid r$, then $\sigma^r=e$ and it is therefore a cycle as well. With that convention the exact condition is
\[\boxed{\gcd(n,r)=1\quad\text{or}\quad n\mid r.}\]
If ``cycle'' is intended to mean a nonidentity $n$-cycle, the familiar condition is just $\gcd(n,r)=1$ (for $n>1$). The extra identity case must not be silently omitted under Definition 9.6.

\exercise{9.36}{Left multiplication is a permutation}
The inverse of $\lambda_a:g\mapsto ag$ is $\lambda_{a^{-1}}$, because $a^{-1}(ag)=g$ and $a(a^{-1}g)=g$. Equivalently, cancellation gives injectivity, and a desired target $h$ is reached by the input $a^{-1}h$, giving surjectivity. Thus $\lambda_a\in S_G$.

\exercise{9.37}{The subgroup of left translations}
By the preceding exercise every $\lambda_a$ is a permutation. The identity is $\lambda_e$. Associativity in $G$ gives
\[(\lambda_a\lambda_b)(g)=a(bg)=(ab)g=\lambda_{ab}(g),\]
so the set is closed under composition. Finally $\lambda_a^{-1}=\lambda_{a^{-1}}$ also lies in it. These facts prove it is a subgroup of $S_G$.

\exercise{9.38}{Left translations act transitively}
Given $x,y\in G$, solve $ax=y$ by taking $a=yx^{-1}$. The left translation $\lambda_a$ then sends $x$ to $y$. Since this works for every pair of points in the set $G$, the left-translation subgroup is transitive on $G$.

\exercise{9.39}{Two generators for the symmetric group}
Put $c=(1,2,\ldots,n)$ and $t=(1,2)$. Conjugating a transposition relabels its entries, so
\[c^k t c^{-k}=(c^k(1),c^k(2)).\]
For $k=0,\ldots,n-2$, these include all adjacent transpositions $s_i=(i,i+1)$, $1\le i<n$. For $i<j$ one has
\[(i,j)=s_i s_{i+1}\cdots s_{j-2}s_{j-1}s_{j-2}\cdots s_{i+1}s_i.\]
To verify it, the product before the middle transposition conjugates $(j-1,j)$ to $(i,j)$; alternatively follow $i,j$ and each intermediate point. Thus every transposition lies in $\gen{c,t}$. Every permutation is a product of transpositions, so $\gen{c,t}=S_n$.
